% 第 3 课时：二次函数与方程、不等式
\chapter{二次函数与方程、不等式}

\begin{objective}
\begin{itemize}
  \item 会用判别式判断方程根的个数，会用求根公式解方程；
  \item 能说出二次函数、二次方程、二次不等式三者的关系；
  \item 会用图象法和数轴穿根法解一元二次不等式；
  \item 会用数轴穿根法解分式不等式。
\end{itemize}
\end{objective}

\begin{keypoint}
重点：一元二次不等式的解法。\\
难点：用图象的正负解释不等式的解集；含参不等式的分类讨论。
\end{keypoint}

\begin{procedure}
\section{温故}
用上一课的配方法，把 $y=2x^2-4x+1$ 写成顶点式，并说出顶点坐标。

\section{判别式与求根公式}
对于方程 $ax^2+bx+c=0$（$a\neq0$），记
\[
  \Delta=b^2-4ac.
\]
\begin{center}
\footnotesize
\setlength{\tabcolsep}{3pt}
\renewcommand{\arraystretch}{1.4}
\begin{tabular}{p{2.8cm}p{4.0cm}p{5.2cm}}
\toprule
$\Delta$ 的符号 & 方程的根 & 函数图象与 $x$ 轴的交点个数 \\
\midrule
$\Delta>0$ & 两个不相等的实数根 & 两个交点 \\
$\Delta=0$ & 两个相等的实数根 & 一个交点（相切） \\
$\Delta<0$ & 没有实数根 & 没有交点 \\
\bottomrule
\end{tabular}
\end{center}

求根公式：
\[
  x=\frac{-b\pm\sqrt{b^2-4ac}}{2a}.
\]

\section{三个二次的关系}
\begin{itemize}
  \item 令 $y=0$，得到一元二次方程 $ax^2+bx+c=0$，它的解就是函数图象与 $x$ 轴交点的横坐标；
  \item 令 $y>0$（或 $y<0$），得到一元二次不等式 $ax^2+bx+c>0$（或 $<0$），
  它的解集就是函数图象在 $x$ 轴上方（或下方）部分所对应的横坐标 $x$ 的取值范围。
\end{itemize}

\begin{figure}[htbp]
  \centering
  \includegraphics[width=0.80\textwidth]{parabola.pdf}
  \caption{$y=x^2-2x-3$：交点、上方、下方对应三种问题}
\end{figure}

\key{一句话记忆}：方程找交点，不等式找上下，根是边界。

\section{一元二次不等式的解法}
\begin{example}
解不等式 $x^2-2x-3<0$。
\end{example}
\begin{solution}
方程 $x^2-2x-3=0$ 的两根为 $x_1=-1$，$x_2=3$。
因为 $a=1>0$，抛物线开口向上，在 $x$ 轴下方的部分是两根之间，
故解集为 $(-1,3)$。
\end{solution}

\begin{example}
解不等式 $-x^2+2x+3>0$。
\end{example}
\begin{solution}
两边同乘 $-1$，不等号方向改变，得 $x^2-2x-3<0$，
由上一例知解集为 $(-1,3)$。
\end{solution}

\key{口诀}：大于取两边，小于取中间（适用于 $a>0$ 的情形）。

\section{数轴穿根法}

对高次不等式，用数轴穿根法更快。它的关键是\key{先看最高次项系数的符号}：
系数为正从右上方起笔；系数为负就先两边同乘 $-1$（不等号方向改变），化为正。
\key{一律先化为正}，规则只有一条。

\begin{example}
解不等式 $(x+1)(x-2)(x-4)<0$。
\end{example}
\begin{solution}
最高次项系数为正。根为 $-1$、$2$、$4$。在数轴上标出后，从右上方起笔依次穿根，
得到各段符号依次为 $+$、$-$、$+$、$-$。
取小于 $0$ 的部分，解集为 $(-\infty,-1)\cup(2,4)$。
\end{solution}

\begin{figure}[htbp]
  \centering
  \includegraphics[width=0.92\textwidth]{ineq-on-axis.pdf}
  \caption{穿根折线与解集（空心点表示端点不取）}
\end{figure}

\key{步骤}：化标准形（最高次项系数化为正）$\to$ 求根标轴 $\to$ 右上起笔 $\to$ 奇穿偶不穿 $\to$ 取上下方 $\to$ 查端点。

\section{分式不等式的解法}

分式不等式也用这条思路：先判断分子、分母的正负，再转化为整式不等式画图求解。
\[
  \frac{f(x)}{g(x)}>0 \iff f(x)g(x)>0 \quad (g(x)\neq0).
\]
\key{注意}：分母的零点取不到，必须画空心点。

\begin{example}
解不等式 $\dfrac{x+1}{x-2}\geq0$。
\end{example}
\begin{solution}
等价于 $(x+1)(x-2)\geq0$ 且 $x\neq2$。
根为 $-1$、$2$，从右上方起笔依次穿根，取大于等于 $0$ 的部分并去掉 $2$，
解集为 $(-\infty,-1]\cup(2,+\infty)$。
\end{solution}

\begin{figure}[htbp]
  \centering
  \includegraphics[width=0.92\textwidth]{rational-ineq.pdf}
  \caption{$\dfrac{x+1}{x-2}\geq0$：$-1$ 实心（可取），$2$ 空心（分母零点）}
\end{figure}

\begin{exercise}
解下列不等式：
\begin{enumerate}
  \item $x^2-5x+6>0$；
  \item $x^2-4x+4\leq0$；
  \item $x^2+2x+3>0$；
  \item $(x-1)(x-2)(x-3)\geq0$。
\end{enumerate}
\end{exercise}
\end{procedure}

\begin{board}
\begin{itemize}
  \item 判别式 $\Delta=b^2-4ac$：判根个数、判交点个数；
  \item 三个二次：令 $y=0$ 得方程，令 $y>0$ 得不等式；
  \item 解不等式：先化正，再求根，后定区间，最后查端点。
\end{itemize}
\end{board}

\begin{homework}
\begin{enumerate}
  \item 解不等式 $x^2-3x-4\geq0$；
  \item 解不等式 $-2x^2+4x-3>0$；
  \item 用数轴穿根法解不等式 $(x+2)(x-1)(x-3)<0$。
\end{enumerate}
\end{homework}
